\documentclass[12pt]{article}
\usepackage{times} 			% use Times New Roman font

\usepackage[margin=1in]{geometry}   % sets 1 inch margins on all sides
\usepackage[hidelinks]{hyperref}               % for URL formatting
\usepackage[pdftex]{graphicx}       % So includegraphics will work

% for fancy page headings
\usepackage{fancyhdr}
\setlength{\headheight}{13.6pt} % to remove fancyhdr warning
\pagestyle{fancy}
\fancyhf{}
\rhead{\small \thepage}
\lhead{\small CS 800, Spring 2026 - A4 - Chitikela}

% more packages and settings
\usepackage{datetime}
\usepackage{indentfirst}  % always indent first paragraph in a section
\usepackage[sort&compress, square, comma, numbers]{natbib}
\usepackage[small, bf]{caption}
\setlength{\parskip}{0.5em}

%-------------------------------------------------------------------------
\begin{document}

\begin{centering}
{\large\textbf{A4 - Literature Review}}\\
\vspace{3mm}
Yamini Chitikela\\
CS 800, Spring 2026\\
\end{centering}

%-------------------------------------------------------------------------

\section{Introduction}

Graph neural networks (GNNs) have emerged as the dominant paradigm for machine
learning on graph-structured data, with applications spanning molecular property
prediction, drug discovery, and social network analysis. Two theoretical
properties lie at the heart of their study: \textit{expressiveness}, the ability
to distinguish structurally non-isomorphic graphs, and \textit{generalization},
the capacity to perform well on unseen graph distributions. This literature review
examines five papers that address both properties through the unifying
mathematical framework of graph homomorphism---edge-preserving mappings between
graphs---which has proven to be a natural and precise language for characterizing
the structural information that GNNs can encode. Barcel\'{o} et al.~\cite{barcelo2021graph}
established both the practical efficiency and theoretical grounding of augmenting
GNNs with homomorphism counts of small graph patterns. Zhang et al.~\cite{zhang2024beyond}
provided a complete quantitative theory of GNN expressiveness by identifying
exactly which graph patterns each architecture can count under homomorphism.
Li et al.~\cite{li2024generalization} connected this structural lens to GNN generalization
through information-theoretic bounds derived from homomorphism entropy.
Li et al.~\cite{li2025towards} directly examined the relationship between expressivity and
generalization, challenging a widely held theoretical assumption. Finally,
Zhang et al.~\cite{zhang2026reinterpreting} extended the homomorphism framework to hypergraph
kernels. These papers were selected because they form a coherent, foundational
body of work that spans from landmark theoretical contributions to the current
research frontier, establishing graph homomorphism as the central analytical tool
for understanding and improving graph learning methods.

\section{Summaries}

Barcel\'{o} et al.~\cite{barcelo2021graph} address the fundamental limitation that
message-passing neural networks (MPNNs) are bounded in expressive power by the
one-dimensional Weisfeiler-Lehman (1-WL) graph isomorphism test, preventing
them from detecting substructures such as triangles or cycles. Their proposed
solution is the $\mathcal{F}$-MPNN framework, which extends any MPNN by
prepending a single preprocessing step that computes homomorphism counts of small
rooted graph patterns from a set $\mathcal{F}$ and appends these counts to each
vertex's initial feature vector. This augmentation maintains the $\mathcal{O}(n)$
per-iteration memory cost of the underlying MPNN. The authors precisely
characterize the expressive power of $\mathcal{F}$-MPNNs through a generalized
algorithm called $\mathcal{F}$-WL: two vertices are indistinguishable by any
$\mathcal{F}$-MPNN if and only if their homomorphism counts from all
$\mathcal{F}$-pattern trees coincide, placing $\mathcal{F}$-MPNNs strictly
between 1-WL and 2-WL in expressiveness. Theoretical results further identify
conditions under which a pattern in $\mathcal{F}$ is redundant given others,
providing practical guidance for pattern selection. Empirical evaluations across
multiple benchmark datasets confirm that augmenting existing GNN architectures
with homomorphism counts consistently yields performance improvements with
negligible computational overhead.

Zhang et al.~\cite{zhang2024beyond} argue that the Weisfeiler-Lehman hierarchy, while
widely used, provides only a coarse and qualitative assessment of GNN
expressiveness that fails to characterize the practical capabilities of modern
architectures. They introduce \textit{homomorphism expressivity} as a complete,
quantitative alternative: for any GNN model $M$, the family $\mathcal{F}^M$
consists of all graph patterns that $M$ can count under graph homomorphism, and
this family uniquely determines the representations that $M$ computes on any
graph. Two GNN models compute identical representations on all graphs if and only
if their homomorphism expressivities coincide, making $\mathcal{F}^M$ a true
fingerprint for each model. Using novel extensions of the nested ear decomposition
(NED) concept from graph theory, the authors derive exact characterizations for
four major GNN families: MPNNs capture forests, Subgraph GNNs capture graphs with
endpoint-shared NED, Local 2-GNNs capture graphs with strong NED, and 2-FGNNs
capture graphs with any NED. These results enable direct quantitative comparison
of architectures, settle several previously open questions, and are empirically
validated by demonstrating that homomorphism expressivity correlates well with
practical model performance.

Li et al.~\cite{li2024generalization} observe that existing generalization bounds for
GNNs typically rely on coarse graph statistics such as maximum degree and tend to
be vacuous for real-world datasets, failing to reflect the structural complexity
of actual graph distributions. Their approach connects GNN generalization to the
\textit{entropy of graph homomorphism}: an information-theoretic quantity that
measures the diversity of homomorphism count vectors across graphs drawn from a
given distribution for a pattern set $\mathcal{F}$. Datasets whose graphs exhibit
low-entropy homomorphism count distributions contain repetitive structural
patterns, which support tighter generalization bounds. The authors derive
non-vacuous generalization bounds for both graph-level and node-level
classification tasks within the $\mathcal{F}$-MPNN framework, covering 1-WL
GNNs, $k$-WL GNNs, Homomorphism-Injected GNNs, and Subgraph-Injected GNNs under
a unified theoretical lens. Empirical comparisons on both real-world benchmark and
synthetic datasets show that the proposed bounds closely track observed
generalization gaps, establishing homomorphism entropy as a practically
meaningful and data-dependent measure of GNN generalization capacity.

Li et al.~\cite{li2025towards} challenge the common theoretical belief that highly
expressive GNNs tend to overfit graph structures and therefore generalize poorly
compared to simpler models. They demonstrate empirically that expressive models
frequently achieve strong generalization performance, and develop a theoretical
framework that reconciles this observation with statistical learning theory. Their
central contribution is a $k$-variance margin-based generalization bound in which
a GNN's generalization gap is linked to the variance in graph structure that the
model captures, as measured relative to a more expressive reference encoder. A key
consequence is that the generalization bound of any GNN can be upper-bounded using
a more expressive encoder, allowing structural complexity to be analyzed through
simpler combinatorial descriptors including homomorphism counts. Their analysis is
architecture-agnostic. They further identify two geometric properties of learned
embedding spaces that govern generalization quality: intra-class concentration,
which measures how tightly same-class graph embeddings cluster, and inter-class
separation in the Wasserstein metric. Experiments on the PROTEINS benchmark
confirm that more expressive models produce embedding distributions exhibiting
both properties simultaneously, supporting the theoretical conclusions.

Zhang et al.~\cite{zhang2026reinterpreting} extend the homomorphism-based expressiveness
framework from ordinary graphs to hypergraph kernels, where hyperedges may
connect any number of vertices to capture higher-order relationships in
bioinformatics and network analysis. They propose a general comparison framework
grounded in hypergraph homomorphism: the expressive power of any hypergraph
kernel is fully characterized by its \textit{pattern space}, the family of
hypergraph substructures whose homomorphism counts the kernel implicitly encodes.
Two kernels are equally expressive if and only if their pattern spaces coincide,
enabling direct comparison without case-by-case analysis. Applying this
framework, the authors characterize the pattern space of the Hypergraph
Weisfeiler-Lehman (HG WL) kernel---the set of all hypergraph subtrees---and
analyze the Hypergraph Rooted kernel, deriving exact conditions under which each
fails to distinguish non-isomorphic hypergraphs. To address these limitations,
they introduce the \textit{Hypergraph Subtree-Cycle Kernel} (HG SCKernel), which
augments tree-based structural features with cycle-based homomorphism count
patterns, yielding two variants of increasing expressiveness. Experiments across
fifteen classification benchmarks confirm the superiority of the proposed kernels
over classical hypergraph kernel methods.

\section{Synthesis}

The five papers reviewed here form a coherent research arc in which graph
homomorphism emerges as the central theoretical tool for analyzing graph learning
methods across two dimensions: expressiveness and generalization.
Barcel\'{o} et al.~\cite{barcelo2021graph} initiated this line of work by showing that homomorphism
counts of small graph patterns can practically and efficiently enhance GNN
expressiveness, establishing the $\mathcal{F}$-MPNN framework as both a
deployable architecture and a theoretical model. This posed the natural question
of exactly what different architectures can compute---a question answered
rigorously and completely by Zhang et al.~\cite{zhang2024beyond}, who proved that the family
of patterns a GNN can count under homomorphism uniquely determines its output
representations. The nested ear decomposition characterizations derived in that
work transformed the study of GNN expressiveness from a qualitative ordering into
a precise, quantitative science: rather than asking whether one model is more
expressive than another, researchers can now identify which specific substructures
separate two architectures in terms of their distinguishing capability.

The subsequent papers extend this expressiveness-oriented framework to
generalization, engaging with the full scope of the learning problem.
Li et al.~\cite{li2024generalization} establish that the information-theoretic entropy of
homomorphism count distributions provides non-vacuous and data-dependent
generalization bounds, directly linking the structural diversity of a graph
dataset to the expected test performance of a GNN. This reframes generalization
analysis from a worst-case exercise dependent on crude graph statistics to a
structurally meaningful one. Li et al.~\cite{li2025towards} complement this by refuting
the assumed expressivity--generalization trade-off, showing theoretically and
empirically that expressive models tend to produce well-concentrated and
well-separated embedding distributions---precisely the geometric properties that
support generalization under margin-based analysis. Finally,
Zhang et al.~\cite{zhang2026reinterpreting} demonstrate that the pattern-space
characterization of expressiveness transfers naturally to hypergraph kernels,
confirming the broad applicability of homomorphism-based analysis beyond standard
GNN architectures. Taken together, these five works establish that structural
information---measured, compared, and designed through the lens of graph
homomorphism---is the fundamental currency of graph learning theory and practice.

%-------------------------------------------------------------------------
\bibliographystyle{ieeetr}
\bibliography{litreview}

\end{document}
